% GENERATED FILE. Edit canonical lessons, then run npm run generate:books.
% canonical-source-hash: fnv1a64:7864ecc4eadade6b
% canonical-lessons: IT-C273-dapprima, IT-C273-innanzitutto, IT-C273-in-seguito, IT-C273-successivamente, IT-C273-dopodiche

\chapter{Prima di tutto}
\label{ch:it-prima-di-tutto}
\hlchaptermodality{273}

\begin{chapteropening}
\textbf{By the end of this chapter:} \emph{I can put the first steps of a story in order: at first, first of all, later on, subsequently and after that.}

Complete the last lesson of chapter 273: I can put the first steps of a story in order: at first, first of all, later on, subsequently and after that.
\end{chapteropening}

\section[dapprima]{dapprima — at first, initially}

\label{lesson:IT-C273-dapprima}



\begin{quote}
Before the new one: say the Italian for exactly, then the Italian for almost.
\end{quote}

\subsection*{You'll want to know: dapprima}
\textbf{dapprima} — \textquotedblleft{}at first, initially\textquotedblright{}.

\begin{quote}
Dapprima non ho capito niente, poi ho riso.
\end{quote}

\begin{quote}
\emph{At first I did not understand anything, then I laughed.}
\end{quote}

A careful, written word; in speech people often say \emph{all'inizio}.

\begin{cousinweb}
\emph{dapprima} is built from \emph{da} \textquotedblleft{}from\textquotedblright{} and \emph{prima} \textquotedblleft{}first, before\textquotedblright{}, which goes back to Latin \emph{prīmus} \textquotedblleft{}first\textquotedblright{} (Lewis \& Short).
\end{cousinweb}

\begin{grammarlens}[title={What you've built}]
One more word for telling, explaining and describing at B1.
\end{grammarlens}

\subsection*{Your turn}
\begin{itemize}
\raggedright
  \item \emph{Say it:} \emph{dapprima}
  \item \emph{Say it:} \emph{dapprima}, once more
  \item \emph{Say it:} a sentence of your own with \emph{dapprima}
  \item \emph{Recall:} say the Italian for exactly, then the Italian for almost, then say \emph{dapprima} again
\end{itemize}

\begin{tcolorbox}[breakable,colback=teal!4,colframe=teal!35!black,title={Before you move on}]
What does dapprima mean? (\textquotedblleft{}At first, initially\textquotedblright{}.) Say it once more.
\end{tcolorbox}
\section[innanzitutto]{innanzitutto — first of all, above all}

\label{lesson:IT-C273-innanzitutto}



\begin{quote}
Before the new one: say the Italian for almost, then the Italian for at first.
\end{quote}

\subsection*{You'll want to know: innanzitutto}
\textbf{innanzitutto} — \textquotedblleft{}first of all, above all\textquotedblright{}.

\begin{quote}
Innanzitutto, grazie a tutti di essere qui.
\end{quote}

\begin{quote}
\emph{First of all, thank you all for being here.}
\end{quote}

Also written \emph{innanzi tutto}; \emph{innanzi} is an old word for ‘before, ahead’.

\begin{cousinweb}
\emph{innanzitutto} joins \emph{innanzi} \textquotedblleft{}before, ahead\textquotedblright{}, usually traced to Late Latin \emph{in ante}, and \emph{tutto} \textquotedblleft{}all\textquotedblright{}: before everything else. Latin \emph{ante} meant \textquotedblleft{}before\textquotedblright{} (Lewis \& Short).
\end{cousinweb}

\begin{grammarlens}[title={What you've built}]
One more word for telling, explaining and describing at B1.
\end{grammarlens}

\subsection*{Your turn}
\begin{itemize}
\raggedright
  \item \emph{Say it:} \emph{innanzitutto}
  \item \emph{Say it:} \emph{innanzitutto}, once more
  \item \emph{Say it:} a sentence of your own with \emph{innanzitutto}
  \item \emph{Recall:} say the Italian for almost, then the Italian for at first, then say \emph{innanzitutto} again
\end{itemize}

\begin{tcolorbox}[breakable,colback=teal!4,colframe=teal!35!black,title={Before you move on}]
What does innanzitutto mean? (\textquotedblleft{}First of all, above all\textquotedblright{}.) Say it once more.
\end{tcolorbox}
\section[in seguito]{in seguito — later on, afterwards}

\label{lesson:IT-C273-in-seguito}



\begin{quote}
Before the new one: say the Italian for at first, then the Italian for first of all.
\end{quote}

\subsection*{You'll want to know: in seguito}
\textbf{in seguito} — \textquotedblleft{}later on, afterwards\textquotedblright{}.

\begin{quote}
In seguito ho trovato un lavoro migliore.
\end{quote}

\begin{quote}
\emph{Later on I found a better job.}
\end{quote}

\emph{In seguito a} + a noun means ‘following, as a result of’: \emph{in seguito all'incidente}.

\begin{cousinweb}
\emph{seguito} is the past participle of \emph{seguire} \textquotedblleft{}to follow\textquotedblright{}, which is inherited from Latin \emph{sequī} \textquotedblleft{}to follow\textquotedblright{} (Lewis \& Short).
\end{cousinweb}

\begin{grammarlens}[title={What you've built}]
One more word for telling, explaining and describing at B1.
\end{grammarlens}

\subsection*{Your turn}
\begin{itemize}
\raggedright
  \item \emph{Say it:} \emph{in seguito}
  \item \emph{Say it:} \emph{in seguito}, once more
  \item \emph{Say it:} a sentence of your own with \emph{in seguito}
  \item \emph{Recall:} say the Italian for at first, then the Italian for first of all, then say \emph{in seguito} again
\end{itemize}

\begin{tcolorbox}[breakable,colback=teal!4,colframe=teal!35!black,title={Before you move on}]
What does in seguito mean? (\textquotedblleft{}Later on, afterwards\textquotedblright{}.) Say it once more.
\end{tcolorbox}
\section[successivamente]{successivamente — subsequently, after that (formal)}

\label{lesson:IT-C273-successivamente}



\begin{quote}
Before the new one: say the Italian for first of all, then the Italian for later on.
\end{quote}

\subsection*{You'll want to know: successivamente}
\textbf{successivamente} — \textquotedblleft{}subsequently, after that (formal)\textquotedblright{}.

\begin{quote}
Successivamente il treno è ripartito.
\end{quote}

\begin{quote}
\emph{Subsequently the train set off again.}
\end{quote}

Formal and common in reports and on the news; in everyday talk, \emph{poi}.

\begin{cousinweb}
\emph{successivamente} is built on \emph{successivo} \textquotedblleft{}following\textquotedblright{}, a learned word usually traced to Late Latin \emph{successīvus}, from \emph{succēdere} \textquotedblleft{}to come after\textquotedblright{} (Lewis \& Short).
\end{cousinweb}

\begin{grammarlens}[title={What you've built}]
One more word for telling, explaining and describing at B1.
\end{grammarlens}

\subsection*{Your turn}
\begin{itemize}
\raggedright
  \item \emph{Say it:} \emph{successivamente}
  \item \emph{Say it:} \emph{successivamente}, once more
  \item \emph{Say it:} a sentence of your own with \emph{successivamente}
  \item \emph{Recall:} say the Italian for first of all, then the Italian for later on, then say \emph{successivamente} again
\end{itemize}

\begin{tcolorbox}[breakable,colback=teal!4,colframe=teal!35!black,title={Before you move on}]
What does successivamente mean? (\textquotedblleft{}Subsequently, after that (formal)\textquotedblright{}.) Say it once more.
\end{tcolorbox}
\section[dopodiché]{dopodiché — after which, and after that}

\label{lesson:IT-C273-dopodiche}



\begin{quote}
Before the new one: say the Italian for later on, then the Italian for subsequently.
\end{quote}

\subsection*{You'll want to know: dopodiché}
\textbf{dopodiché} — \textquotedblleft{}after which, and after that\textquotedblright{}.

\begin{quote}
Abbiamo cenato, dopodiché siamo usciti.
\end{quote}

\begin{quote}
\emph{We had dinner, after which we went out.}
\end{quote}

Joins two events in a row. Also written \emph{dopo di che}.

\begin{cousinweb}
\emph{dopodiché} fuses \emph{dopo di che} \textquotedblleft{}after which\textquotedblright{}. \emph{Dopo} is usually traced to Late Latin \emph{dē post}, built on Latin \emph{post} \textquotedblleft{}after\textquotedblright{} (Lewis \& Short).
\end{cousinweb}

\begin{grammarlens}[title={What you've built}]
One more word for telling, explaining and describing at B1.
\end{grammarlens}

\subsection*{Your turn}
\begin{itemize}
\raggedright
  \item \emph{Say it:} \emph{dopodiché}
  \item \emph{Say it:} \emph{dopodiché}, once more
  \item \emph{Say it:} a sentence of your own with \emph{dopodiché}
  \item \emph{Recall:} say the Italian for later on, then the Italian for subsequently, then say \emph{dopodiché} again
\end{itemize}

\begin{tcolorbox}[breakable,colback=teal!4,colframe=teal!35!black,title={Before you move on}]
What does dopodiché mean? (\textquotedblleft{}After which, and after that\textquotedblright{}.) Say it once more.
\end{tcolorbox}
